\section{问题重述}

无创产前检测（NIPT）通过分析孕妇外周血中的游离胎儿DNA片段来评估胎儿的遗传状态。当胎儿为男性时，母体血液中可检测到Y染色体特异性序列，其浓度随孕周增长而升高。孕妇的体质量指数（BMI）越高，母体自身游离DNA对胎儿DNA的稀释效应越强，导致检测信噪比下降。

本文围绕以下四个问题展开建模。问题一要求建立Y染色体浓度关于孕周与BMI的定量模型，刻画浓度随孕周的非线性增长趋势及BMI的稀释作用。问题二要求按BMI水平对孕妇分组，在每组内确定NIPT的最优检测时点，使检测的准确性与时效性达到综合最优。问题三将模型扩展至包含年龄、体重、身高等多因素，评估不同体征组合下的检测达标率，并分析测量误差在模型中的传播效应。问题四针对女性胎儿（无Y染色体信号），利用Z值等指标建立染色体非整倍体异常的判别模型。
